\section{Proposed Method}
\subsection{Issue-time Data Contract}
For station $s$ and target day $d$, $X_{s,d}$ is assembled before issue time,
whereas $y_{s,d}$ is constructed later from measured pollutant concentrations.
Observations may contribute only when their bins end by 17:00 IST on $d-1$;
completed daily-history inputs are from $d-2$ or earlier. Weather summaries
use archived forecasts with a documented 48-hour Previous Runs offset.
Every admitted snapshot passes the availability guard. Historical observation
publication times are unavailable, so local-history skill is conditional on
the protocol's healthy-feed latency assumption.

\begin{figure}[t]\centering\footnotesize\AirGuardArchitecture
\caption{AirGuard study architecture. The primary model comparisons and
selection lock remain pending; prospective forecasts are a separate immutable
stream. Targets are not used to choose issue-time predictors.}\label{fig:architecture}
\end{figure}

\subsection{Measured Targets and Sensor Quality}
OpenAQ preserves the CPCB-provider station and sensor provenance
\cite{openaq2026}. Values must have convertible mass units, positive finite
concentrations, at least 75\% provider coverage and no provider rejection flag.
Exact zeros and identical sequences of six or more consecutive bins are
rejected under v3.1. Potential spikes and PM2.5 above PM10 remain flagged
observations. Feature QC uses the prefix visible at the cutoff; later readings
cannot revise an earlier feature. Rejected station-days remain in the dataset.

The common CPCB-style engine uses at least 16 distinct hours for each supported
pollutant, at least three supported species and a required particulate species
\cite{cpcb2014aqi}. Particulates and designated gases use 24-hour means;
CO and O$_3$ use the maximum complete trailing eight-hour mean, including
seven preceding support hours. For a bounded breakpoint interval,
\begin{equation}
I_p(C)=I_L+\frac{I_H-I_L}{C_H-C_L}(C-C_L),\qquad
I=\max_{p\in S}I_p(C_p).\label{eq:aqi}
\end{equation}
The result is an AirGuard CPCB-style index, not an official CPCB bulletin.
The Severe category is open-ended: $y\geq401$ is a lower bound.
Such targets are excluded from exact regression fitting and MAE/RMSE/bias;
category, exceedance, below-401 and lower-bound shortfall diagnostics retain them.

\subsection{Operational CAMS Reconstruction}
The ADS dataset is the global atmospheric-composition forecast archive,
not reanalysis \cite{cams2024forecasts}. GRIB type, producing date and time,
lead hours, valid time, grid and units are decoded and verified. The downloaded
grid is 0.4$^\circ$; matching uses the frozen nearest grid centre. At the declared
issue time the study uses the $d-1$ 00:00 UTC run, with nominal dissemination
treated explicitly as an assumption rather than an empirical delivery log.

Surface PM mass concentrations are converted from kg m$^{-3}$ to $\mu$g m$^{-3}$.
Gas mass mixing ratio $x$ uses the same run's lowest-model-level temperature
$T$, specific humidity $q$ and hybrid pressure $p$:
\begin{equation}
\rho=\frac{p}{287.05T(1+0.6078q)},\qquad C_{\mu\mathrm{g/m^3}}=10^9\rho x.
\label{eq:density}
\end{equation}
CO uses $10^6$ for mg m$^{-3}$. Model-level fields are interpolated within their
producing run before conversion. Missing brackets, hybrid coefficients or
inconsistent metadata cause rejection. Retrieval time records acquisition
provenance; it does not establish historical publication time.

\subsection{Direct and Residual Calibration}
Let $b_{s,d}$ denote the frozen raw CAMS AQI. Direct learners predict $y_{s,d}$;
residual learners fit exact measured-minus-CAMS labels:
\begin{equation}
r_{s,d}=y_{s,d}-b_{s,d},\qquad
\widehat y_{s,d}=\max\{0,b_{s,d}+f_\theta(X_{s,d})\}.\label{eq:residual}
\end{equation}
Ridge uses standardised, training-imputed features with an L2 penalty.
LightGBM and XGBoost provide nonlinear alternatives. Missing-value statistics
and learned parameters are fitted only on training rows. Predictor groups
comprise CAMS concentrations and descriptors, measured history, forecast
weather, calendar information and spatial context. Fire features remain
optional and require admissible publication evidence before evaluation.
Target-specific pollutant support is an evaluation annotation, never a predictor.
The current valid targets use the frozen PM2.5/PM10/O$_3$ support; the
six-pollutant index is a secondary comparison representing the live product.

\subsection{Uncertainty and Advisories}
Split conformal calibration uses exact calibration residuals $e_i$ with
\begin{equation}
\begin{aligned}
k&=\lceil(n+1)(1-\alpha)\rceil,\quad q_\alpha=e_{(k)},\\
\mathcal C(x)&=[\max(0,\widehat y-q_\alpha),\widehat y+q_\alpha].
\end{aligned}
\label{eq:conformal}
\end{equation}
If the finite-sample rank cannot be supported, an interval is unavailable.
Conformalised quantile LightGBM and forecast-category Mondrian calibration
are additional declared methods. Membership and quantile-crossing repair use
forecasts, never observed test categories. Censored-Severe compatibility
($U\geq401$) is not identified coverage of the unknown true index.

Advisory policies use the point forecast (P0), calibrated median (P1), fixed
0.75 quantile (P2), or the smallest candidate quantile meeting a calibration-only
under-warning target (P3). For observed category $c$ and warning category $w$,
\begin{equation}
L_-=\mathbf1\{w<c\},\quad L_+=\mathbf1\{w>c\},\quad
\widehat R_-=\frac{k_-+1}{n+1}.\label{eq:warning}
\end{equation}
The declared targets are 10\% for general adults and 5\% for the other exposure
strata. No feasible candidate is reported as such. Under dependence, these
are empirical goals rather than a guaranteed personal or clinical risk bound.

\subsection{Feature Availability Audit}
The feature contract distinguishes a value's valid time from when it could
have been available. A forecast can describe future valid hours while still
being an admissible pre-issue input; an observation cannot. Retrieving an old
field today is insufficient without its producing-run identity. The same
distinction applies to archived weather, measured-history latency and fire
publication. Table~\ref{tab:availability} documents the admitted groups and
remaining assumptions.
\begin{table}[t]\centering\footnotesize
\caption{Feature availability and operational assumptions. Static metadata and the calendar do not require outcome information.}
\label{tab:availability}
\setlength{\tabcolsep}{4pt}
\begin{tabular}{lp{.62\columnwidth}}
\toprule
Group & Admitted information\\\midrule
A CAMS & Verified operational run, concentrations, sub-indices and diurnal descriptors\\
B History & Bins ending by 17:00 IST; completed daily lags from $d-2$; healthy-feed assumption\\
C Weather & Archived previous-run forecasts with pre-issue run bound\\
D Calendar & Target-day month, weekday and cyclical season information\\
E Spatial & Frozen coordinates, terrain and admissible pre-issue neighbours\\
F Fire & Optional; withheld without historical publication evidence\\
\bottomrule\end{tabular}
\end{table}
\subsection{Baseline Support and Missingness}
The primary CAMS comparison does not credit a learner for removing a species
the station failed to measure. Frozen PM2.5/PM10/O$_3$ support is defined from
sensor metadata; the current valid join has no support mismatch. The
six-pollutant index remains a secondary comparator because it is closer to
the live product's display. Both are produced by the same engine, so target
and baseline breakpoints, averaging windows and Severe treatment cannot drift.

History lags use station/date keys rather than row offsets. Missing days remain
missing; neither targets nor forecast indices are borrowed from the wrong date.
Censored historical values are not fabricated exact bias labels. Neighbour
summaries use cutoff-compliant measurements and exclude self-donations.
Monitor-free spatial regimes must remove unavailable donor feeds, not merely
drop a station identifier. Missing predictors are imputed using training-only
statistics with indicators, preserving the conditions under which a forecast
could be issued.

\subsection{Advisory Evaluation Contract}
Under-warning and over-warning are evaluated jointly. Serious under-warning
(two or more categories), misses on Poor-or-worse days and warning distance
are separate diagnostics. P1 and P2 use one-sided calibration of their fixed
quantiles; P3 selects among declared raw quantile forecasts on the calibration
window. That selection is itself its calibration step. No test-window outcome
can select a warning level or repair an interval.

The candidates are $\{0.5,0.65,0.75,0.85,0.9,0.95,0.975\}$.
An absent feasible candidate is a failed empirical risk target, not an excuse
to relabel a permissive warning as controlled. Shared station outcomes across
the seven exposure strata are not seven independent clinical samples.
The final report must retain unfavourable trade-offs, including cases where
reduced under-warning requires substantial over-warning.
